\chapter{Structureur, ventes et sales-trader}\label{ch:in:structurer-sales-and-sales-trader}

L'association européenne des émetteurs de produits structurés publie chaque trimestre ce que détiennent et
négocient les marchés de ses membres. Fin mars 2026, les investisseurs de six marchés européens détenaient 487 milliards d'euros de produits d'investissement structurés
émis sous forme de titres, principalement en Suisse et en Allemagne, et 163\,413 nouveaux produits d'investissement ont été cotés au
seul premier trimestre. Derrière chaque note se trouvent un \omterm{def:in:structurer-sales-and-sales-trader:structurer}{structureur} qui l'a conçue, un vendeur qui l'a placée, un
desk de trading qui la couvre, et une marge que l'acheteur de la note paie et que le desk de la banque encaisse. Ce chapitre décrit les trois métiers
tournés vers la clientèle où la compétence quantitative fait partie du poste, calcule la marge d'une note type,
et expose les règles qui s'imposent à ceux qui vendent des produits financiers.

\begin{center}\footnotesize
\begin{tabular}{@{}p{2.4cm}p{3.3cm}p{3.3cm}p{3.3cm}@{}}
\toprule
\multicolumn{4}{@{}l}{\textbf{Fiches de rôle~: les métiers tournés vers la clientèle}}\\
\midrule
 & \omterm{def:in:structurer-sales-and-sales-trader:structurer}{structureur} & \omterm{def:in:structurer-sales-and-sales-trader:sales}{vendeur institutionnel} & sales-trader\\
\midrule
travaille sur & conception et valorisation des produits & relations avec les clients et idées & ordres des clients et leur exécution\\
horizon & semaines à années (la durée de vie de la note) & mois à années & secondes à jours\\
crédité de & la marge du desk & chiffre d'affaires client & flux client\\
codes dans les déclarations & 13-2099.01, 13-2051, 41-3031 & 41-3031 & 41-3031, 13-2099.01\\
enseigné dans & livre~5, ch.~18--19 & livre~9, ch.~24 & livre~1, ch.~2~; livre~10, ch.~20\\
\bottomrule
\end{tabular}
\end{center}

\section{Le structureur}

\begin{definition}[Structureur]\label{def:in:structurer-sales-and-sales-trader:structurer}
Un \emph{structureur}\index{structureur} est un spécialiste quantitatif de la salle des marchés d'une banque qui conçoit des produits structurés pour les clients
et les distributeurs~: il choisit le profil de gain, le sous-jacent, la protection et le coupon, les valorise avec les modèles du desk,
et fixe la marge de structuration dans le cadre de l'approbation du produit et de la capacité du desk de trading à le couvrir.
\end{definition}

Un \omterm{def:in:structurer-sales-and-sales-trader:structurer}{structureur} travaille entre trois parties. Le client, ou la banque privée ou le réseau de détail qui vend aux clients, veut
un coupon ou un niveau de protection~; le desk de trading qui couvrira la note veut des risques qu'il peut porter (livre~5, chapitre~27~;
livre~9, chapitre~28)~; la banque veut une marge. Le livre~5, chapitre~19, décrit l'activité, et le chapitre~18 traite de
l'autocallable, le profil de gain que ce chapitre évalue. Le travail quantitatif du \omterm{def:in:structurer-sales-and-sales-trader:structurer}{structureur} consiste à valoriser le produit et
ses sensibilités, à comprendre de quels paramètres dépend le prix et lesquels le client ne voit pas.

\begin{definition}[Gouvernance des produits]\label{def:in:structurer-sales-and-sales-trader:governance}
La \emph{gouvernance des produits}\index{gouvernance des produits} désigne l'ensemble des processus par lesquels une société qui fabrique ou
distribue des instruments financiers approuve chaque produit pour un marché cible identifié de clients finaux, vérifie que la
stratégie de distribution convient à ce marché, et réexamine le produit tout au long de sa vie.
\end{definition}

Le produit du \omterm{def:in:structurer-sales-and-sales-trader:structurer}{structureur} doit passer par le processus d'approbation des produits de la société avant d'être vendu~: les règles européennes exigent qu'un
fabricant «~maintienne, fasse fonctionner et réexamine un processus d'approbation de chaque instrument financier~» qui précise
«~un marché cible identifié de clients finaux~». Le marché cible est une déclaration sur la personne qui devrait acheter la note~; le
\omterm{def:in:structurer-sales-and-sales-trader:structurer}{structureur} en rédige une grande partie.

\section{La marge dans la note}

\begin{method}[Marge de structuration d'une note]\label{met:in:structurer-sales-and-sales-trader:margin}
Une note vendue au pair, 100, a une fair value $V(c)$ pour 100 selon le modèle du desk pour le coupon $c$~; sa marge est $100-V(c)$.
Pour un autocallable Phoenix, les flux de la note sont linéaires en fonction du coupon, donc le long des mêmes trajectoires simulées $V(c)=V(0)+c\,\partial V/\partial c$,
et le coupon qui laisse une marge cible $m$ est
$c_m=\bigl(100-m-V(0)\bigr)/(\partial V/\partial c)$.
\end{method}

\begin{example}[Un autocallable de cinq ans]\label{ex:in:structurer-sales-and-sales-trader:note}
Une note d'illustration sur un indice actions~: cinq ans, constatations annuelles, rappel anticipé lorsque l'indice est égal ou supérieur à son
niveau initial, un coupon de 6 pour 100 par an versé lorsque l'indice est égal ou supérieur à 70~\% (avec mémoire), capital remboursé sauf si
l'indice termine en dessous de 60~\% de son niveau initial, auquel cas le porteur subit la perte. Rendement de dividende 2~\%, volatilité plate (flat),
200\,000 trajectoires. À 6~\%, la marge est de $-1.65$ pour 100 avec des taux à 1~\% et une volatilité de 20~\% (la banque perdrait), 1.83
avec des taux à 4~\% et une volatilité de 20~\%, 6.91 à 1~\% et 30~\%, et 9.13 à 4~\% et 30~\%. Le coupon qui laisse une marge de 2~\%
vaut 4.23~\% à 1~\% et 20~\%, 5.91~\% à 4~\% et 20~\%, 8.79~\% à 1~\% et 30~\%, et 10.22~\% à 4~\% et 30~\%
(\cref{fig:in:structurer-sales-and-sales-trader:coupon}).
\end{example}

\begin{center}\footnotesize
\begin{tabular}{@{}lrrrr@{}}
\toprule
 taux~; volatilité & 1~\%~; 20~\% & 4~\%~; 20~\% & 1~\%~; 30~\% & 4~\%~; 30~\%\\
\midrule
marge à un coupon de 6~\%, pour 100 & $-1.65$ & 1.83 & 6.91 & 9.13\\
coupon pour une marge de 2~\%, \% par an & 4.23 & 5.91 & 8.79 & 10.22\\
coupon pour une marge nulle, \% par an & 5.20 & 6.97 & 9.92 & 11.41\\
rappel anticipé à la première année, \% & 44.0 & 50.0 & 42.7 & 46.7\\
perte de capital à l'échéance, \% & 15.8 & 9.4 & 25.1 & 19.3\\
\bottomrule
\end{tabular}
\end{center}

Deux enseignements constituent le matériel quotidien du \omterm{def:in:structurer-sales-and-sales-trader:structurer}{structureur}. Des taux plus élevés permettent d'obtenir, pour une même marge, un coupon plus élevé~: la note est
une obligation plus des options vendues par l'investisseur, et une obligation actualisée à 4~\% coûte moins cher qu'à 1~\%. Une volatilité plus élevée permet
également un coupon plus élevé, car l'investisseur est vendeur de la baisse et en est rémunéré~; mais la probabilité de perdre du capital augmente
avec elle, de 15.8~\% à 25.1~\% des trajectoires à des taux de 1~\%. Un client qui compare les coupons d'une note à l'autre compare,
sans le voir, la volatilité vendue.

\begin{omfigure}
\begin{tikzpicture}
\begin{axis}[ybar,area legend,width=10.5cm,height=5.4cm,xmin=0.5,xmax=4.5,ymin=0,ymax=13,bar width=14pt,xtick={1,2,3,4},
  xticklabels={{1~\%, 20~\%},{4~\%, 20~\%},{1~\%, 30~\%},{4~\%, 30~\%}},xlabel={\small taux, volatilité},
  ylabel={\small coupon, \% par an},tick label style={font=\footnotesize},grid=major,grid style={gray!20},
  table/col sep=comma,legend style={font=\scriptsize,at={(0.03,0.97)},anchor=north west,legend cell align=left}]
\addplot[fill=omDef,draw=none] table[x=pos,y=fair]{figdata/industry/23-structurer-sales-and-sales-trader/coupon.csv};
\addplot[fill=omDef!35,draw=none] table[x=pos,y=zero]{figdata/industry/23-structurer-sales-and-sales-trader/coupon.csv};
\legend{coupon laissant une marge de 2~\%,coupon laissant une marge nulle}
\end{axis}
\end{tikzpicture}

\omcaption{Le coupon de l'autocallable d'illustration sur cinq ans du chapitre qui laisse une marge de 2~\% et une marge nulle, à deux
niveaux de taux et deux volatilités. L'écart entre chaque paire est la marge exprimée en coupon. Données~:
\texttt{in\_structurer.margins}, issues de \texttt{firm.roles.structuring\_margin} sur \texttt{firm.autocall}.}\label{fig:in:structurer-sales-and-sales-trader:coupon}
\end{omfigure}

\section{Ventes institutionnelles et couverture}

\begin{definition}[Vendeur institutionnel]\label{def:in:structurer-sales-and-sales-trader:sales}
Un \emph{vendeur institutionnel}\index{vendeur institutionnel} est une personne qui couvre un ensemble de
clients professionnels (sociétés de gestion, hedge funds, fonds de pension, assureurs, entreprises) pour une banque ou un courtier~: elle leur apporte les
idées, les prix et les produits de la société, porte leur intérêt auprès des desks, et est créditée du chiffre d'affaires que génèrent les clients.
\end{definition}

Le vendeur est propriétaire de la relation et de la franchise que formalise la segmentation des clients (tiering) du livre~9, chapitre~24. Lorsque
les marchés se négocient par voie électronique, la dimension quantitative du métier est l'analyse d'exécution du client, la tarification du desk et
le risque qu'une transaction laisse à la banque~; un vendeur capable de discuter des trois est plus utile au client. Le rôle est
jugé sur le chiffre d'affaires du client et sur la part de son activité que la banque capte, mesures que la banque suit et auxquelles la rémunération du vendeur est indexée
(chapitre~13).

\section{Le sales-trader}

Le sales-trader du livre~1, chapitre~2, se place entre les ordres des clients et le marché~: il prend l'ordre d'un \omterm{def:in:portfolio-manager-and-pod-analyst:roles}{gérant de portefeuille},
décide avec le client de la manière de le traiter, le fait passer par les algorithmes de la société ou par ses traders, et rend compte du résultat
par rapport à l'indice de référence du client. Le trading high-touch du livre~10, chapitre~20, correspond à ce métier. Son contenu quantitatif est
l'exécution~: analyse des coûts de transaction, choix de l'algorithme et de la plateforme, impact de marché des gros ordres (livre~10), et
le sales-trader capable d'expliquer une exécution par rapport au prix d'arrivée conserve le client.

\section{Compétences quantitatives dans les métiers en contact avec la clientèle}

Les trois rôles mobilisent la compétence quantitative de manière différente.
\begin{itemize}
\item \textbf{Le \omterm{def:in:structurer-sales-and-sales-trader:structurer}{structureur}} tarifie et conçoit~: calcul stochastique et Monte Carlo (livre~5), sensibilités d'un
payoff à la volatilité, à la corrélation et aux taux, et coût de couverture que le desk devra supporter.
\item \textbf{Le vendeur} traduit~: transforme la vision d'un desk ou un produit en problème pour le client, et inversement, et a besoin
d'assez de modèle pour savoir quand un prix est faux.
\item \textbf{Le sales-trader} mesure~: qualité d'exécution par rapport aux indices de référence, en chiffres que le client peut vérifier.
\end{itemize}
Dans les trois cas, le public est un client qui peut en savoir moins, et la compétence consiste aussi à dire ce que le client achète.

\section{Les règles applicables au travail en contact avec la clientèle}

Les personnes qui vendent des produits financiers sont agréées et supervisées selon des modalités que les quants d'une société de trading ne connaissent pas.

\begin{dated}{2026-09}{Règles applicables aux personnes qui conçoivent et vendent des produits}\label{dat:in:structurer-sales-and-sales-trader:rules}
Union européenne~: MiFID II (directive 2014/65/UE) article 16(3) exige un processus d'approbation des produits par le producteur, avec «~un
marché cible identifié de clients finaux~»~; l'article 24(2) exige des produits «~conçus pour répondre aux besoins d'un
marché cible identifié~» et une stratégie de distribution compatible~; l'article 25(1) exige que le personnel qui conseille ou informe les clients
«~possède les connaissances et la compétence nécessaires~». Les orientations de l'ESMA sur la \omterm{def:in:structurer-sales-and-sales-trader:governance}{gouvernance des produits} au titre de MiFID II (ESMA35-43-3448, 3
août 2023) précisent comment les articles 16(3) et 24(2) s'appliquent. États-Unis~: un représentant général en valeurs mobilières réussit
l'examen Series 7 de la FINRA, qui «~évalue la compétence d'un représentant enregistré débutant~», la Securities
Industry Essentials exam étant exigée en complément~; la Series 63, examen de la NASAA administré par la FINRA, porte sur le droit des valeurs mobilières des États.
\end{dated}

Pour un quant qui passe à un rôle en contact avec la clientèle, ces règles modifient le métier de trois façons. Le produit doit être approuvé
avant d'être vendu et revu pendant toute sa durée de vie~; la personne doit être qualifiée pour l'activité~; et les registres de la société
incluent «~l'enregistrement des conversations téléphoniques ou des communications électroniques~» relatives aux ordres des clients (MiFID II,
article 16(7)). Le chapitre~24 décrit la fonction de conformité qui les fait respecter.

\begin{dated}{2026-03}{Le marché des produits structurés}\label{dat:in:structurer-sales-and-sales-trader:market}
Rapport de l'EUSIPA pour le premier trimestre 2026~: ventes de produits d'investissement et à effet de levier sur les plateformes de négociation dans les
marchés déclarants, 74 milliards d'euros (produits d'investissement 19 milliards)~; encours de produits d'investissement structurés à la fin de mars,
486\,635 millions d'euros (Suisse 286\,264 millions, Allemagne 99\,890, Italie 66\,788, Autriche
17\,462, Belgique 11\,678, Luxembourg 4\,553)~; 465\,721 produits d'investissement cotés~; 163\,413 lancés au cours du trimestre.
\end{dated}

\begin{omfigure}
\begin{tikzpicture}
\begin{axis}[xbar,width=10.5cm,height=4.8cm,xmin=0,xmax=330,ymin=0.4,ymax=6.6,bar width=9pt,ytick=data,
  yticklabels={{Luxembourg},{Belgique},{Autriche},{Italie},{Allemagne},{Suisse}},
  yticklabel style={font=\footnotesize},xticklabel style={font=\footnotesize},
  xlabel={\small encours de produits d'investissement structurés, fin mars 2026, milliards d'euros},grid=major,grid style={gray!20},
  table/col sep=comma,nodes near coords={\pgfplotspointmeta\%},nodes near coords style={font=\scriptsize,anchor=west,
  xshift=1pt},
  point meta=explicit symbolic]
\addplot[fill=omDef,draw=none] table[x=bn,y=pos,meta=share]{figdata/industry/23-structurer-sales-and-sales-trader/eusipa.csv};
\end{axis}
\end{tikzpicture}

\omcaption{Encours de produits d'investissement structurés émis sous forme de titres dans six marchés européens à la fin de mars 2026,
rapport de l'EUSIPA~; les étiquettes donnent la part de chaque marché dans le total. Données~: rapport de marché EUSIPA T1 2026,
via \texttt{in\_structurer.eusipa\_shares}.}\label{fig:in:structurer-sales-and-sales-trader:eusipa}
\end{omfigure}

Le marché est concentré~: la Suisse détient 58.8~\% de l'encours des six marchés et l'Allemagne 20.5~\%
(\cref{fig:in:structurer-sales-and-sales-trader:eusipa}). Un desk qui émet entre un et cinq milliards d'euros de titres par an
à une marge de 2~\% dégage 20 à 100 millions d'euros de chiffre d'affaires de structuration~; réparti sur dix \omterm{def:in:structurer-sales-and-sales-trader:structurer}{structureurs}, 2 à 10 millions d'euros
chacun. L'émission comme l'équipe sont illustratives~; l'arithmétique montre pourquoi la rémunération d'un \omterm{def:in:structurer-sales-and-sales-trader:structurer}{structureur} suit la marge
et le volume que le desk peut placer.

\begin{dated}{2025-05}{Ce que gagnent les agents commerciaux~: l'enquête}\label{dat:in:structurer-sales-and-sales-trader:pay}
Enquête sur les professions, mai 2025, agents commerciaux en valeurs mobilières, matières premières et services financiers (41-3031), qui comprend
des \omterm{def:in:structurer-sales-and-sales-trader:sales}{vendeurs institutionnels}, des sales-traders et des courtiers~: secteur des valeurs mobilières 173\,040 salariés, médiane 103\,030~\$, 10\up{e}~centile à
90\up{e}~centile~: 55\,130 à 309\,440~\$~; banques 257\,110 salariés, médiane 61\,440~\$. Les intitulés de \omterm{def:in:structurer-sales-and-sales-trader:structurer}{structureur} et de vente des déclarations sont trop
peu nombreux pour être publiés séparément. Salaires seulement.
\end{dated}

\begin{omfigure}
\begin{tikzpicture}
\begin{axis}[width=10.5cm,height=4.4cm,xmin=0,xmax=330,ymin=0.4,ymax=3.6,ytick=data,
  yticklabels={{banques},{finance et assurance},{valeurs mobilières}},
  yticklabel style={font=\footnotesize},xticklabel style={font=\footnotesize},xlabel={\small salaire annuel, milliers de \$},
  grid=major,grid style={gray!20},table/col sep=comma]
\addplot[draw=none,mark=none,error bars/.cd,x dir=both,x explicit,error mark=none,error bar style={omDef,line width=1pt}]
  table[x=p50,y=pos,x error plus expr=\thisrow{p90}-\thisrow{p50},x error minus expr=\thisrow{p50}-\thisrow{p10}]
  {figdata/industry/23-structurer-sales-and-sales-trader/survey.csv};
\addplot[draw=none,mark=none,error bars/.cd,x dir=both,x explicit,error mark=none,error bar style={omDef!40,line width=7pt}]
  table[x=p50,y=pos,x error plus expr=\thisrow{p75}-\thisrow{p50},x error minus expr=\thisrow{p50}-\thisrow{p25}]
  {figdata/industry/23-structurer-sales-and-sales-trader/survey.csv};
\addplot[only marks,mark=|,mark size=5pt,line width=1.5pt,omThm] table[x=p50,y=pos]{figdata/industry/23-structurer-sales-and-sales-trader/survey.csv};
\end{axis}
\end{tikzpicture}

\omcaption{Salaires des agents commerciaux en valeurs mobilières, matières premières et services financiers (41-3031) par secteur, enquête
sur les professions de mai 2025~: 10\up{e}~centile à 90\up{e}~centile (fin), écart interquartile (épais), médiane (repère). La profession comprend
aussi les courtiers de détail en plus des ventes institutionnelles. Données~: \texttt{data/industry/oews\_roles.csv}, via
\texttt{in\_structurer.survey}.}\label{fig:in:structurer-sales-and-sales-trader:survey}
\end{omfigure}

\section{Tutoriel~: la marge de la note}\label{tut:in:structurer-sales-and-sales-trader:tut}

\begin{tutorial}
\textbf{Objectif.} Tarifer un autocallable, trouver sa marge et le coupon qui laisse une marge cible, et mettre le chiffre
d'affaires du desk par \omterm{def:in:structurer-sales-and-sales-trader:structurer}{structureur} en regard de la taille du marché.
\textbf{État final~:} \cref{fig:in:structurer-sales-and-sales-trader:coupon,fig:in:structurer-sales-and-sales-trader:eusipa}
et le tableau des marges.
\begin{tutsteps}
\item \textbf{Fiche de termes.} \texttt{firm.autocall.TermSheet} du livre~5, avec des observations annuelles, un coupon Phoenix avec
mémoire et protection à l'échéance.
\item \textbf{Valeur et marge.} \texttt{firm.roles.note\_value} simule l'indice et actualise les flux~;
\texttt{structuring\_margin} utilise la linéarité de la note en coupon pour trouver le coupon de n'importe quelle marge
(\cref{lst:in:structurer-sales-and-sales-trader:margin}).
\omcode{code/firm/roles/firm_roles.py}{408}{416}{La marge de structuration et le coupon qui laisse une marge cible.}\label{lst:in:structurer-sales-and-sales-trader:margin}
\item \textbf{Grille.} Deux niveaux de taux et deux volatilités, les mêmes trajectoires pour chaque coupon.
\item \textbf{Échelle.} Chiffre d'affaires par \omterm{def:in:structurer-sales-and-sales-trader:structurer}{structureur} à partir d'une fourchette d'émissions et d'une équipe illustratives.
\end{tutsteps}
\end{tutorial}

Pour la grille de l'exemple, le tableau ci-dessus donne le résultat~; au coupon de 6~\%, faire passer les taux de 1~\% à 4~\% ajoute 3.48 points de
marge à une volatilité de 20~\%.

\textbf{Pour aller plus loin.} Ajouter le spread de financement de l'émetteur, que paie aussi l'acheteur de la note~; valoriser
sous volatilité locale (livre~5, chapitre~18) et comparer~; faire de la note un worst-of sur trois indices et observer le coupon
et la probabilité de perte monter ensemble.

\section{Construction~: cartes orientées client et marge de structuration}\label{bld:in:structurer-sales-and-sales-trader:roles}

\begin{build}
\bfield{Rôle} Ajouter les rôles orientés client au registre et calculer la marge dans une note.
\bfield{Interface} \texttt{firm.roles}~: les fiches \texttt{structurer}, \texttt{institutional salesperson},
\texttt{sales-trader}~; \texttt{note\_value(sheet, r, q, vol, n, seed, steps)}~;
\texttt{structuring\_margin(make\_sheet, coupon, target, r, q, vol, n, seed, steps)}~; sur \texttt{firm.autocall}.
\bfield{Règles} Valeurs pour 100 vendus au pair~; les mêmes nombres aléatoires pour chaque coupon, de sorte que le coupon pour une marge soit
exact le long des trajectoires~; volatilité plate sauf si l'appelant passe une fonction.
\bfield{Tests d'acceptation} \texttt{code/firm/roles/tests/}~: la note réévaluée au coupon renvoyé a exactement la marge cible~;
le coupon pour une marge nulle est supérieur au coupon pour une marge positive.
\bfield{Pour aller plus loin} Notes worst-of~; la sensibilité de la marge à chaque paramètre~; un générateur de fiches de termes qui cherche le
niveau de protection pour un objectif de coupon.
\end{build}

\begin{omsources}
\item EUSIPA, Market Report on Structured Investment and Leverage Products Q1/2026.
\item Directive 2014/65/EU (MiFID II), Articles 16, 24 et 25~; ESMA35-43-3448, Guidelines on MiFID II product governance.
\item FINRA, Series 7 and Series 63 qualification exams.
\item Bureau of Labor Statistics, occupational survey, May 2025~; livre~5, chapitres~18 et~19.
\end{omsources}

\section{Exercices}

\begin{exercise}[$\star$]\label{exo:in:structurer-sales-and-sales-trader:1}
Quelle est la part des produits structurés en circulation sur les six marchés que la Suisse et
l'Allemagne détenaient ensemble à la fin de mars 2026~?
\end{exercise}

\begin{exercise}[$\star$]\label{exo:in:structurer-sales-and-sales-trader:2}
À des taux de 4~\% et une volatilité de 20~\%, quelle est la marge de la note à 6~\%, et quel coupon laisserait une marge nulle~?
\end{exercise}

\begin{exercise}[$\star$]\label{exo:in:structurer-sales-and-sales-trader:3}
Quel article de MiFID II impose un processus d'approbation des produits par le producteur, et que doit-il préciser~?
\end{exercise}

\begin{exercise}[$\star\star$]\label{exo:in:structurer-sales-and-sales-trader:4}
Pourquoi un taux d'intérêt plus élevé permet-il un coupon plus élevé pour une même marge~?
\end{exercise}

\begin{exercise}[$\star\star$]\label{exo:in:structurer-sales-and-sales-trader:5}
Un client préfère la note à 10.22~\% (taux 4~\%, volatilité 30~\%) à la note à 5.91~\% (taux 4~\%, volatilité 20~\%).
Qu'achète-t-il avec le coupon supplémentaire~?
\end{exercise}

\begin{exercise}[$\star\star$]\label{exo:in:structurer-sales-and-sales-trader:6}
Pourquoi la rémunération d'un \omterm{def:in:structurer-sales-and-sales-trader:structurer}{structureur} est-elle liée à la marge et au volume du desk plutôt qu'à la performance ultérieure de la note~?
\end{exercise}

\begin{exercise}[$\star\star\star$]\label{exo:in:structurer-sales-and-sales-trader:7}
\emph{Programmation.} Valorisez la note à 6~\% à des taux de 4~\% et une volatilité de 20~\% avec 50\,000 trajectoires et les graines de 1 à 20.
Quel est l'écart type de la marge selon les graines, et le signe de la marge est-il incertain~?
\end{exercise}

\begin{exercise}[$\star\star\star$]\label{exo:in:structurer-sales-and-sales-trader:8}
\emph{Trouver la faille.} «~La note verse 10~\% par an et protège le capital jusqu'à 60~\%, donc elle est plus sûre que de détenir
l'indice et rapporte davantage que ses dividendes.~»
\end{exercise}

\section{Problème~: la marge dans la note}

\begin{problem}[{Problème du week-end --- la marge dans la note}]\label{pb:in:structurer-sales-and-sales-trader:1}
Un analyste quantitatif se voit proposer un poste dans un desk de structuration et veut comprendre d'où vient son chiffre d'affaires
et quelles sont les règles du métier.

\medskip\noindent\textbf{Partie I --- Les rôles.}
\begin{enumerate}
\item Définir le \omterm{def:in:structurer-sales-and-sales-trader:structurer}{structureur}, le \omterm{def:in:structurer-sales-and-sales-trader:sales}{vendeur institutionnel} et la \omterm{def:in:structurer-sales-and-sales-trader:governance}{gouvernance des produits}.
\item Entre qui un \omterm{def:in:structurer-sales-and-sales-trader:structurer}{structureur} travaille-t-il~?
\item Que fait un sales-trader, et qu'est-ce que le trading high-touch~?
\item Quelle compétence quantitative chaque rôle mobilise-t-il~?
\item Quel livre enseigne l'autocallable~?
\end{enumerate}

\medskip\noindent\textbf{Partie II --- La marge.}
\begin{enumerate}[resume]
\item Énoncer la note de l'exemple.
\item Définir la marge de structuration et en déduire le coupon pour une marge cible.
\item Donner la marge à 6~\% dans les quatre cas.
\item Donner le coupon pour une marge de 2~\% dans les quatre cas.
\item Pourquoi la volatilité augmente-t-elle le coupon, et que fait-elle d'autre~?
\end{enumerate}

\medskip\noindent\textbf{Partie III --- Le marché et les règles.}
\begin{enumerate}[resume]
\item Donner la taille du marché et sa concentration par pays.
\item Quel chiffre d'affaires par \omterm{def:in:structurer-sales-and-sales-trader:structurer}{structureur} le desk de l'exemple réalise-t-il~?
\item Énoncer les règles de \omterm{def:in:structurer-sales-and-sales-trader:governance}{gouvernance des produits} et de compétence de MiFID II.
\item Que couvrent les Series 7 et 63 de la FINRA~?
\item Que montrent les salaires de l'enquête~?
\end{enumerate}

\medskip\noindent\textbf{Partie IV --- Le verdict.}
\begin{enumerate}[resume]
\item Énoncer le \emph{résultat nommé}~: la marge de structuration de l'autocallable du chapitre à des taux de 1~\% et 4~\%, et le
chiffre d'affaires par \omterm{def:in:structurer-sales-and-sales-trader:structurer}{structureur} pour une émission du desk dans la fourchette du chapitre.
\item Que changerait un spread de financement~?
\item Que doit un \omterm{def:in:structurer-sales-and-sales-trader:structurer}{structureur} au client final~?
\item Qu'apporteraient les compétences quantitatives de l'analyste au desk, et que n'apporteraient-elles pas~?
\item En deux phrases, conseiller l'analyste.
\end{enumerate}
\end{problem}

\section{Questions d'entretien}

\begin{interviewq}[$\star$ \iqroles{banque}]\label{iq:in:structurer-sales-and-sales-trader:1}
Décomposez une note à capital protégé en une obligation et une option.
\end{interviewq}

\begin{interviewq}[$\star$ \iqroles{banque, trader}]\label{iq:in:structurer-sales-and-sales-trader:2}
Un client demande pourquoi l'autocallable de cette année rapporte davantage que celui de l'an dernier. Donnez trois raisons.
\end{interviewq}

\begin{interviewq}[$\star\star$ \iqroles{banque}]\label{iq:in:structurer-sales-and-sales-trader:3}
L'investisseur dans un autocallable est-il long ou court en volatilité, et long ou court en corrélation dans un worst-of~?
\end{interviewq}

\begin{interviewq}[$\star\star$ \iqroles{banque, trader}]\label{iq:in:structurer-sales-and-sales-trader:4}
L'ordre de participation de 5~\% d'un client s'est terminé 20 points de base moins bien que le prix d'arrivée. Comment l'expliquez-vous~?
\end{interviewq}

\begin{interviewq}[$\star\star$ \iqroles{banque, risque}]\label{iq:in:structurer-sales-and-sales-trader:5}
Qu'est-ce qui entre dans une déclaration de marché cible pour une note comme celle du chapitre~?
\end{interviewq}

\begin{interviewq}[$\star\star\star$ \iqroles{banque, chercheur}]\label{iq:in:structurer-sales-and-sales-trader:6}
Le desk a vendu de nombreux autocallables sur le même indice. Quel risque le desk porte-t-il, et quand fait-il mal~?
\end{interviewq}
